% Generato da tools/notion2latex.py a partire dalla pagina Notion "ripasso-probabilita".

\chapter{Ripasso di probabilità di base}

\section{1. Le basi}

\textbf{Spazio campionario} \(\Omega\): l'insieme di tutti i risultati possibili di un esperimento casuale.

\textbf{Evento} \(A\): un sottoinsieme di \(\Omega\), cioè un gruppo di risultati che ci interessa.

\begin{esempio}

Esempio: lanci un dado. \(\Omega = \{1,2,3,4,5,6\}\). L'evento ``esce pari'' è \(A = \{2,4,6\}\).

\end{esempio}

\subsection{Assiomi}

\begin{itemize}
\tightlist
\item
  \(0 \le \Pr[A] \le 1\): una probabilità sta sempre tra 0 (impossibile) e 1 (certo)
\item
  \(\Pr[\Omega] = 1\): qualcosa succede di sicuro
\item
  se \(A\) e \(B\) sono \textbf{disgiunti} (non possono succedere insieme): \(\Pr[A \cup B] = \Pr[A] + \Pr[B]\)
\end{itemize}

\begin{esempio}

Esempio: dado. ``Esce 1'' e ``esce 2'' sono disgiunti, quindi \(\Pr[\text{1 o 2}] = 1/6 + 1/6 = 1/3\).

\end{esempio}

\subsection{Distribuzione uniforme}

Tutti i risultati hanno la stessa probabilità. Allora:

\[
\Pr[A] = \frac{|A|}{|\Omega|} = \frac{\text{casi favorevoli}}{\text{casi possibili}}
\]

\begin{esempio}

Esempio: pesco una carta da un mazzo da 52. \(\Pr[\text{asso}] = 4/52 = 1/13\).

\end{esempio}

\subsection{Complementare}

\[
\Pr[\neg A] = 1 - \Pr[A]
\]

Utilissimo quando ``almeno una volta'' è difficile da contare ma ``mai'' è facile.

\begin{esempio}

Esempio: lanci un dado 4 volte. Probabilità di almeno un 6?

Contare direttamente è scomodo (un 6, due 6, \ldots). Il complementare è ``nessun 6'': ogni lancio non dà 6 con prob. \(5/6\), quindi

\(\Pr[\text{almeno un 6}] = 1 - (5/6)^4 = 1 - 625/1296 \approx 0{,}518\)

\end{esempio}

\subsection{Unione di eventi non disgiunti}

\[
\Pr[A \cup B] = \Pr[A] + \Pr[B] - \Pr[A \cap B]
\]

Si sottrae l'intersezione perché altrimenti la conti due volte.

\begin{esempio}

Esempio: pesco una carta. Probabilità che sia di cuori \textbf{o} una figura (J, Q, K)?

\(\Pr[\text{cuori}] = 13/52\), \(\Pr[\text{figura}] = 12/52\), \(\Pr[\text{figura di cuori}] = 3/52\)

\(\Pr = 13/52 + 12/52 - 3/52 = 22/52 \approx 0{,}42\)

\end{esempio}

\section{2. Union bound}

\[
\Pr[A_1 \cup \dots \cup A_m] \le \sum_{i=1}^{m} \Pr[A_i]
\]

Vale \textbf{sempre}, anche se gli eventi non sono disgiunti: sommare le probabilità al massimo sovrastima (perché conti più volte le intersezioni). Si usa quando il valore esatto è difficile ma basta un limite superiore.

\begin{esempio}

Esempio: hai 3 lampadine, ognuna si brucia entro l'anno con probabilità 1\%. La probabilità che \textbf{almeno una} si bruci è al più \(3 \times 1\% = 3\%\).

Il valore esatto (se sono indipendenti) è \(1 - 0{,}99^3 \approx 2{,}97\%\): il bound è molto vicino.

\end{esempio}

\section{3. Probabilità condizionata e indipendenza}

\subsection{Probabilità condizionata}

Con \(\Pr[B] > 0\):

\[
\Pr[A \mid B] = \frac{\Pr[A \cap B]}{\Pr[B]}
\]

Si legge ``probabilità di \(A\) sapendo che è successo \(B\)''. In pratica: restringi il mondo ai soli casi in cui vale \(B\), e guardi quanti di quelli soddisfano anche \(A\).

\begin{esempio}

Esempio: lanci un dado e un amico ti dice ``è uscito un numero pari''. Probabilità che sia 6?

\(\Pr[6 \mid \text{pari}] = \dfrac{\Pr[6]}{\Pr[\text{pari}]} = \dfrac{1/6}{1/2} = \dfrac{1}{3}\)

Ha senso: i casi possibili sono ora solo \(\{2,4,6\}\), e uno su tre è il 6.

\end{esempio}

\subsection{Indipendenza}

\(A\) e \(B\) sono \textbf{indipendenti} se

\[
\Pr[A \cap B] = \Pr[A] \cdot \Pr[B]
\]

che equivale a \(\Pr[A \mid B] = \Pr[A]\): sapere che è successo \(B\) non cambia nulla su \(A\).

\begin{esempio}

Esempio sì: lanci due monete. ``Prima testa'' e ``seconda testa'' sono indipendenti: \(\Pr[\text{TT}] = 1/2 \cdot 1/2 = 1/4\).

Esempio no: peschi due carte \textbf{senza rimetterle}. \(\Pr[\text{1ª asso}] = 4/52\), ma \(\Pr[\text{2ª asso} \mid \text{1ª asso}] = 3/51\): la prima pescata cambia la seconda.

Esempio sorprendente: dado, \(A\) = ``pari'' \(= \{2,4,6\}\), \(B\) = ``al massimo 4'' \(= \{1,2,3,4\}\).

\(\Pr[A \cap B] = \Pr[\{2,4\}] = 2/6 = 1/3\) e \(\Pr[A]\Pr[B] = 1/2 \cdot 2/3 = 1/3\). Sono indipendenti anche se riguardano lo stesso dado!

\end{esempio}

\textbf{Attenzione}: disgiunti ≠ indipendenti. Due eventi disgiunti con probabilità positiva sono anzi \textbf{dipendenti}: se succede uno, l'altro è impossibile.

\subsection{Probabilità totale}

Se \(B_1, \dots, B_m\) partizionano \(\Omega\) (coprono tutti i casi e non si sovrappongono):

\[
\Pr[A] = \sum_{i} \Pr[A \mid B_i]\,\Pr[B_i]
\]

Cioè: spezzi il problema in casi, calcoli ogni caso e fai la media pesata.

\begin{esempio}

Esempio: hai due urne. Urna 1: 3 palline rosse e 2 blu. Urna 2: 1 rossa e 4 blu. Lanci una moneta per scegliere l'urna, poi peschi una pallina.

\(\Pr[\text{rossa}] = \Pr[\text{rossa} \mid U_1]\Pr[U_1] + \Pr[\text{rossa} \mid U_2]\Pr[U_2] = \tfrac{3}{5}\cdot\tfrac{1}{2} + \tfrac{1}{5}\cdot\tfrac{1}{2} = 0{,}4\)

\end{esempio}

\subsection{Teorema di Bayes}

\[
\Pr[A \mid B] = \frac{\Pr[B \mid A]\,\Pr[A]}{\Pr[B]}
\]

Serve a ``invertire'' la condizionata: conosci \(\Pr[B \mid A]\) ma ti interessa \(\Pr[A \mid B]\).

\begin{esempio}

Esempio (il classico test medico): una malattia colpisce l'1\% della popolazione. Il test è positivo nel 99\% dei malati, ma dà anche un 5\% di falsi positivi sui sani. Fai il test ed è positivo: che probabilità hai di essere malato?

Prima \(\Pr[\text{pos}]\) con la probabilità totale:

\(\Pr[\text{pos}] = 0{,}99 \cdot 0{,}01 + 0{,}05 \cdot 0{,}99 = 0{,}0099 + 0{,}0495 = 0{,}0594\)

Poi Bayes:

\(\Pr[\text{malato} \mid \text{pos}] = \dfrac{0{,}99 \cdot 0{,}01}{0{,}0594} \approx 0{,}167\)

Solo il 17\% circa! Il motivo: i sani sono tantissimi, quindi anche un 5\% di falsi positivi su di loro supera di molto i veri positivi.

\end{esempio}

\section{4. Variabili aleatorie e valore atteso}

\subsection{Variabile aleatoria}

Una \textbf{variabile aleatoria} \(X\) è una funzione \(\Omega \to \mathbb{R}\): associa un numero a ogni risultato dell'esperimento.

\begin{esempio}

Esempio: lanci due dadi, \(X\) = somma dei due. Se esce (3, 5), \(X = 8\). \(\Pr[X = 2] = 1/36\) (solo (1,1)), \(\Pr[X = 7] = 6/36\) (sei combinazioni).

\end{esempio}

\subsection{Valore atteso}

\[
\mathbb{E}[X] = \sum_x x \cdot \Pr[X = x]
\]

È la media dei valori pesata con le loro probabilità: il risultato ``medio'' se ripeti l'esperimento tantissime volte.

\begin{esempio}

Esempio: un dado. \(\mathbb{E}[X] = (1+2+3+4+5+6)/6 = 3{,}5\). Non esce mai 3,5, ma è la media a lungo andare.

Esempio (scommessa): paghi 1 € per giocare, se esce 6 vinci 5 €. Il guadagno netto \(X\) vale \(+4\) con prob. \(1/6\) e \(-1\) con prob. \(5/6\).

\(\mathbb{E}[X] = 4 \cdot \tfrac{1}{6} - 1 \cdot \tfrac{5}{6} = -\tfrac{1}{6} \approx -0{,}17\) €. In media perdi 17 centesimi a giocata.

\end{esempio}

\subsection{Linearità del valore atteso}

\[
\mathbb{E}[X + Y] = \mathbb{E}[X] + \mathbb{E}[Y] \qquad \mathbb{E}[aX] = a\,\mathbb{E}[X]
\]

Vale \textbf{sempre}, anche se \(X\) e \(Y\) non sono indipendenti.

\begin{esempio}

Esempio: somma di due dadi. Invece di fare la tabella 6×6: \(\mathbb{E}[X_1 + X_2] = 3{,}5 + 3{,}5 = 7\).

\end{esempio}

\subsection{Variabili indicatrici}

\(X_A = 1\) se l'evento \(A\) succede, \(0\) altrimenti. Allora \(\mathbb{E}[X_A] = \Pr[A]\).

Combinate con la linearità servono a \textbf{contare} quante volte succede qualcosa in media: scrivi il conteggio come somma di indicatrici.

\begin{esempio}

Esempio: lanci 10 monete. Numero medio di teste? \(X = X_1 + \dots + X_{10}\) dove \(X_i = 1\) se la moneta \(i\) dà testa. \(\mathbb{E}[X] = 10 \cdot 1/2 = 5\).

Esempio (guardaroba): \(n\) persone lasciano il cappello, che vengono restituiti a caso. Quante persone in media riavranno il proprio?

\(X_i = 1\) se la persona \(i\) riprende il suo: \(\Pr = 1/n\). Quindi \(\mathbb{E}[X] = n \cdot 1/n = 1\), \textbf{qualunque sia} \(n\). Le \(X_i\) non sono indipendenti, ma la linearità funziona lo stesso.

\end{esempio}

\subsection{Variabili indipendenti}

\(X\) e \(Y\) sono indipendenti se \(\Pr[X=x \wedge Y=y] = \Pr[X=x]\Pr[Y=y]\) per ogni \(x, y\). In tal caso vale (senza indipendenza in generale no):

\[
\mathbb{E}[XY] = \mathbb{E}[X]\,\mathbb{E}[Y]
\]

\begin{esempio}

Esempio: prodotto di due dadi lanciati separatamente: \(\mathbb{E}[X_1 X_2] = 3{,}5 \cdot 3{,}5 = 12{,}25\).

Controesempio: \(X\) = un dado, \(Y = X\) (lo stesso dado). \(\mathbb{E}[X \cdot X] = (1+4+9+16+25+36)/6 \approx 15{,}17 \ne 12{,}25\).

\end{esempio}
